% ═══════════════════════════════════════════════════════════
%  ОСНОВНОЙ КЛАСС И ГЕОМЕТРИЯ
% ═══════════════════════════════════════════════════════════
\usepackage[
    top=2.5cm,
    bottom=2.5cm,
    left=2.5cm,
    right=2.5cm,
    headheight=14pt
]{geometry}

% ══════════════════════════════════════════════════════════
%  ЯЗЫК И КОДИРОВКА
% ══════════════════════════════════════════════════════════
% ─── Язык (через polyglossia, не babel!) ───
\usepackage{polyglossia}
\setmainlanguage{russian}
\setotherlanguage{english}
%\usepackage[english,russian]{babel}

% ═════════════════════════════════════════════════════════
%  ШРИФТЫ
% ═════════════════════════════════════════════════════════
% ─── Шрифты (через fontspec) ───
\usepackage{fontspec}

% Основной — CMU Serif (Computer Modern Unicode, с кириллицей!)
\setmainfont{CMU Serif}

% Без засечек — CMU Sans Serif
\setsansfont{CMU Sans Serif}

% Моноширинный — DejaVu Sans Mono (у него гарантированно есть кириллица)
\setmonofont{DejaVu Sans Mono}[Scale=0.85]

% Русские переносы
\usepackage{hyphenat}

% Математика
\usepackage{amsmath, amssymb, amsfonts}

% ─── Красивые подписи рисунков и таблиц ──────────
%\usepackage{caption}
%\captionsetup{
%    labelsep=period,  % Точка после номера
%    font=small        % Чуть мельче
%}

\usepackage{caption}
\captionsetup[listing]{
    labelsep=space,
    name=Листинг,
    justification=raggedright,
    singlelinecheck=false,
}


% ══════════════════════════════════════════════════════════
%  ПОЛЕЗНЫЕ ПАКЕТЫ
%═══════════════════════════════════════════════════════════

% Многострочные комментарии
\usepackage{comment}

% Ссылки в PDF
\usepackage[
    colorlinks=true,
    linkcolor=clrIdea,
    citecolor=clrIdea,
    urlcolor=clrIdea,
    bookmarks=true,
    bookmarksnumbered=true,
    unicode=true,                        % ← ДОБАВИТЬ
    pdftitle={STM32, DMA и GTK4},
    pdfauthor={Александр},
]{hyperref}

% Умные ссылки: \autoref{fig:...} → "рис. 5"
\usepackage{cleveref}

% Иконки
\usepackage{fontawesome5}

% Единицы измерения
\usepackage{siunitx}
\sisetup{
    output-decimal-marker={,},
    group-separator={\,},
    group-minimum-digits=4,
}

% Улучшенная типографика
\usepackage{microtype}

% Листинги программ
\usepackage{minted}
\setminted{autogobble}
% Настройка по умолчанию для всех листингов
\setminted{
    fontsize=\small,
    linenos=true,
    numbersep=8pt,
    frame=single,
    framesep=2mm,
    breaklines=true,
    tabsize=4,
    bgcolor=clrCodeBg,
    rulecolor=clrCodeFrame,
}

% ─── Список листингов: русский язык и формат ──────────────
\renewcommand{\listingscaption}{Листинг}
\renewcommand{\listoflistingscaption}{Список листингов}

% Красивая подсветка под C
% или 'friendly', 'tango', 'xcode'
\usemintedstyle{default}

% Заметки на полях
\usepackage{marginnote}

% Зачёркнутый текст
\usepackage{ulem}

% Многострочные комментарии. Дублирует comment
%\usepackage{verbatim}

% Условная компиляция (черновик / финал)
\usepackage{ifdraft}

% Скруглённые углы у includegraphics
\usepackage[export]{adjustbox}

\usepackage{epigraph}
%\usepackage{imakeidx}
%\makeindex[title=Предметный указатель, intoc=true]

\usepackage{tocloft}

\usepackage{xcolor}

% ─── TikZ и библиотеки ───────────────────────────────────────
\usepackage{tikz}
\usetikzlibrary{arrows.meta, positioning, shapes, shapes.geometric,
                calc, fit, backgrounds, decorations.pathreplacing,
                decorations.markings, shadows, fadings, patterns,
                matrix, chains}

\usepackage{needspace}

% резисторы прямоугольниками
%\usepackage[europeanresistors]{circuitikz}
% полный европейский стиль (и конденсаторы, и катушки)
\usepackage[european]{circuitikz}

% ─── Типографика: растягивание строк ──────────────────────
% LaTeX разрешает растянуть пробелы 
% на 3em, чтобы избежать Overfull
\setlength{\emergencystretch}{3em}
% разрешает более свободные пробелы (по умолчанию 200)
\tolerance=1000
% убирает предупреждения Underfull \hbox
\hbadness=10000
% убирает предупреждения Underfull \vbox
\vbadness=10000

\usepackage{siunitx}
\sisetup{
    output-decimal-marker={,},
    group-separator={\,},
    group-minimum-digits=4,
}

% Регистрируем единицы, которых нет в siunitx по умолчанию
\DeclareSIUnit{\sample}{Sa}      % семпл
\DeclareSIUnit{\bit}{bit}        % бит
\DeclareSIUnit{\bitps}{bit/s}    % бит в секунду
\DeclareSIUnit{\baud}{Bd}        % бод
\DeclareSIUnit{\sps}{sps}        % samples per second (альтернатива)

%══════════════════════════════════════════════════════════
%           ОФОРМЛЕНИЕ ПОДПИСЕЙ 
%  (точка вместо двоеточия, двойная нумерация)
%══════════════════════════════════════════════════════════

%\AtBeginDocument
%{
% эти команды выполнятся после 
% начала документа, когда всё готово.
%}

% ─── Рисунки и таблицы: точка вместо двоеточия ─────
\AtBeginDocument{%
    \renewcommand{\figurename}{Рис.}
    \renewcommand{\thefigure}{\thechapter.\arabic{figure}}

    \renewcommand{\tablename}{Таблица}
    \renewcommand{\thetable}{\thechapter.\arabic{table}}
}

% Этот код даёт фантомную страницу в начале документа
% ─── Листинги: "Листинг 1.1." + нумерация по главе ─
%\AtBeginDocument{%
%\renewcommand{\thelisting}{\thechapter.\arabic{listing}}%
%\renewcommand{\fnum@listing}{\listingscaption~\thelisting.}%
%}

\usepackage{caption}
\captionsetup[listing]{labelsep=period, name=Листинг}
\renewcommand{\thelisting}{\thechapter.\arabic{listing}}

% ─── cleveref: "рис. 1.1", "табл. 1.1", "лист. 1.1" ─
\crefname{figure}{рис.}{рис.}
\Crefname{figure}{Рис.}{Рис.}
\crefname{table}{табл.}{табл.}
\Crefname{table}{Табл.}{Табл.}
\crefname{listing}{лист.}{лист.}
\Crefname{listing}{Лист.}{Лист.}



\endinput

